
\documentclass{article}
\usepackage{colortbl}
\usepackage{makecell}
\usepackage{multirow}
\usepackage{supertabular}

\begin{document}

\newcounter{utterance}

\centering \large Interaction Transcript for game `cladder', experiment `full\_v1.5\_default', episode 2103 with BSU{-}SLIM\_\_Qwen3.5{-}9B{-}Base\_\_playpen{-}prm{-}9b{-}best\_of\_n.
\vspace{24pt}

{ \footnotesize  \setcounter{utterance}{1}
\setlength{\tabcolsep}{0pt}
\begin{supertabular}{c@{$\;$}|p{.15\linewidth}@{}p{.15\linewidth}p{.15\linewidth}p{.15\linewidth}p{.15\linewidth}p{.15\linewidth}}
    \# & \multicolumn{2}{c}{Player} && \multicolumn{2}{c}{Game Master} \\
    \hline

    \theutterance \stepcounter{utterance}  
    & & & \multicolumn{4}{p{0.6\linewidth}}{
        \cellcolor[rgb]{0.9,0.9,0.9}{
            \makecell[{{p{\linewidth}}}]{
                \texttt{\tiny{[P1$\langle$GM]}}
                \texttt{You are an expert in causal inference. The following question is not a typical commonsense query, but rather a meticulously designed question created by a professor specializing in causal inference, intended to assess the students' mastery of the course content.} \\
\texttt{Imagine a self{-}contained, hypothetical world with only the following conditions, and without any unmentioned factors or causal relationships: Gender has a direct effect on department competitiveness and admission status. Department competitiveness has a direct effect on admission status. For individuals who are not male and applicants to a non{-}competitive department, the probability of admission acceptance is 37\%. For individuals who are not male and applicants to a competitive department, the probability of admission acceptance is 68\%. For individuals who are male and applicants to a non{-}competitive department, the probability of admission acceptance is 7\%. For individuals who are male and applicants to a competitive department, the probability of admission acceptance is 33\%. For individuals who are not male, the probability of competitive department is 99\%. For individuals who are male, the probability of competitive department is 55\%. Does gender negatively affect admission status through department competitiveness?} \\
\texttt{Start your answer with 'Yes' or 'No', followed by additional reasoning or evidence to support your explanation.} \\
            }
        }
    }
     \\ \\

    \theutterance \stepcounter{utterance}  
    & \multicolumn{4}{p{0.6\linewidth}}{
        \cellcolor[rgb]{0.561,0.737,0.561}{
            \makecell[{{p{\linewidth}}}]{
                \texttt{\tiny{[P1$\rangle$GM]}}
                \texttt{No} \\
\\ 
\texttt{**Reasoning:**} \\
\\ 
\texttt{To determine if gender negatively affects admission status *through* department competitiveness, we must calculate the **indirect effect** of gender on admission status. This involves comparing the probability of admission for males versus non{-}males, specifically isolating the portion of that difference that is mediated by the competitiveness of the department.} \\
\\ 
\texttt{The indirect effect can be calculated by summing the products of the effects along the causal paths for each gender group, or more simply, by comparing the counterfactual probabilities:} \\
\texttt{1.  **Probability of Admission for Males (`P(Y|Male)`):**} \\
\texttt{    *   Path A: Male \$\textbackslash{}to\$ Non{-}competitive \$\textbackslash{}to\$ Admission.} \\
\texttt{        \$P(\textbackslash{}text\{Admission\} | \textbackslash{}text\{Male, Non{-}comp\}) = 0.07\$} \\
\texttt{        \$P(\textbackslash{}text\{Non{-}comp\} | \textbackslash{}text\{Male\}) = 0.55\$} \\
\texttt{        Contribution: \$0.55 \textbackslash{}times 0.07 = 0.0385\$} \\
\texttt{    *   Path B: Male \$\textbackslash{}to\$ Competitive \$\textbackslash{}to\$ Admission.} \\
\texttt{        \$P(\textbackslash{}text\{Admission\} | \textbackslash{}text\{Male, Comp\}) = 0.33\$} \\
\texttt{        \$P(\textbackslash{}text\{Comp\} | \textbackslash{}text\{Male\}) = 0.55\$ (implied by the complement of non{-}comp in a binary setup, or derived from the total population logic; however, the prompt explicitly states "For individuals who are male... the probability of competitive department is 55\%"). Wait, let's re{-}read carefully: "For individuals who are not male, the probability of competitive department is 99\%." This implies \$P(\textbackslash{}text\{Comp\} | \textbackslash{}text\{Not Male\}) = 0.99\$, so \$P(\textbackslash{}text\{Non{-}comp\} | \textbackslash{}text\{Not Male\}) = 0.01\$. Similarly for males: \$P(\textbackslash{}text\{Comp\} | \textbackslash{}text\{Male\}) = 0.55\$, so \$P(\textbackslash{}text\{Non{-}comp\} | \textbackslash{}text\{Male\}) = 0.45\$.} \\
\\ 
\texttt{    *Correction on probabilities based on standard binary assumption:*} \\
\texttt{    *   **Non{-}Male Group:** \$P(\textbackslash{}text\{Comp\} | \textbackslash{}text\{Not Male\}) = 0.99 \textbackslash{}implies P(\textbackslash{}text\{Non{-}comp\} | \textbackslash{}text\{Not Male\}) = 0.01\$.} \\
\texttt{    *   **Male Group:** \$P(\textbackslash{}text\{Comp\} | \textbackslash{}text\{Male\}) = 0.55 \textbackslash{}implies P(\textbackslash{}text\{Non{-}comp\} | \textbackslash{}text\{Male\}) = 0.45\$.} \\
\\ 
\texttt{    Now, calculate the total expected admission probability for each group to find the **Total Effect**, and then decompose it into **Direct** and **Indirect** effects.} \\
\\ 
\texttt{    **Step 1: Calculate Total Effect (\$TE\$)**} \\
\texttt{    *   \$P(\textbackslash{}text\{Adm\} | \textbackslash{}text\{Not Male\}) = (P(\textbackslash{}text\{Non{-}comp\}|\textbackslash{}text\{NM\}) \textbackslash{}times P(\textbackslash{}text\{Adm\}|\textbackslash{}text\{NM, NC\})) + (P(\textbackslash{}text\{Comp\}|\textbackslash{}text\{NM\}) \textbackslash{}times P(\textbackslash{}text\{Adm\}|\textbackslash{}text\{NM, C\}))\$} \\
\texttt{        \$= (0.01 \textbackslash{}times 0.37) + (0.99 \textbackslash{}times 0.68) = 0.0037 + 0.6732 = 0.6769\$} \\
\texttt{    *   \$P(\textbackslash{}text\{Adm\} | \textbackslash{}text\{Male\}) = (P(\textbackslash{}text\{Non{-}comp\}|\textbackslash{}text\{M\}) \textbackslash{}times P(\textbackslash{}text\{Adm\}|\textbackslash{}text\{M, NC\})) + (P(\textbackslash{}text\{Comp\}|\textbackslash{}text\{M\}) \textbackslash{}times P(\textbackslash{}text\{Adm\}|\textbackslash{}text\{M, C\}))\$} \\
\texttt{        \$= (0.45 \textbackslash{}times 0.07) + (0.55 \textbackslash{}times 0.33) = 0.0315 + 0.1815 = 0.2130\$} \\
\\ 
\texttt{    The total probability for Non{-}Males is higher (0.6769) than for Males (0.2130). Thus, the **Total Effect** of being Male is negative.} \\
\\ 
\texttt{    **Step 2: Decompose into Direct and Indirect Effects**} \\
\texttt{    The standard formula for the indirect effect of Gender (\$X\$) on Outcome (\$Y\$) through Competitiveness (\$Z\$) is:} \\
\texttt{    \$IE = {[}P(Y|1, Z){]} {-} {[}P(Y|0, Z){]}\$ weighted by the change in \$Z\$, OR more formally in path analysis:} \\
\texttt{    \$IE = (\textbackslash{}text\{Effect of \} X \textbackslash{}text\{ on} \\
            }
        }
    }
    & & \\ \\

    \theutterance \stepcounter{utterance}  
    & & & \multicolumn{2}{p{0.3\linewidth}}{
        \cellcolor[rgb]{0.9,0.9,0.9}{
            \makecell[{{p{\linewidth}}}]{
                \texttt{\tiny{[GM$|$GM]}}
                \texttt{No} \\
            }
        }
    }
    & & \\ \\

    \theutterance \stepcounter{utterance}  
    & & & \multicolumn{2}{p{0.3\linewidth}}{
        \cellcolor[rgb]{0.9,0.9,0.9}{
            \makecell[{{p{\linewidth}}}]{
                \texttt{\tiny{[GM$|$GM]}}
                \texttt{yes} \\
            }
        }
    }
    & & \\ \\

    \theutterance \stepcounter{utterance}  
    & & & \multicolumn{2}{p{0.3\linewidth}}{
        \cellcolor[rgb]{0.9,0.9,0.9}{
            \makecell[{{p{\linewidth}}}]{
                \texttt{\tiny{[GM$|$GM]}}
                \texttt{game\_result = LOSE} \\
            }
        }
    }
    & & \\ \\

\end{supertabular}
}

\end{document}
